% Introducción íntegra de EGA I, pp.5--9; no es el tomo completo. Fragmento editable para el proyecto del tomo.

% Unidad R565: introducción, p. 5, dedicatoria y requisitos a)--d).
% Continuación: intro-fr.tex, párrafo «Ce n'est pas le lieu...».
\section*{Introducción}

\begin{center}
\emph{A Oscar Zariski y André Weil.}
\end{center}

\oldpage[I]{5}

Este trabajo, y los numerosos trabajos que han de seguirlo, están destinados a formar un tratado sobre los fundamentos de la geometría algebraica.
En principio no presuponen ningún conocimiento particular de esta disciplina; incluso se ha comprobado que tal conocimiento, pese a sus evidentes ventajas, podía a veces resultar perjudicial para quien desea familiarizarse con el punto de vista y las técnicas expuestos aquí, por el hábito de adoptar de manera demasiado exclusiva el punto de vista birracional que dicho conocimiento lleva consigo.
En cambio, supondremos que el lector conoce bien los temas siguientes:

\begin{enumerate}
  \item[a)] El \emph{álgebra conmutativa}, tal como se expone, por ejemplo, en los volúmenes de los \emph{Éléments} de N.~Bourbaki que están en preparación (y, hasta que aparezcan esos volúmenes, en Samuel-Zariski~\cite{I-13} y Samuel~\cite{I-11},~\cite{I-12}).
  \item[b)] El \emph{álgebra homológica}, para la cual remitimos a Cartan-Eilenberg~\cite{I-2} (citado como~(M)) y Godement~\cite{I-4} (citado como~(G)), así como al reciente artículo de A.~Grothendieck~\cite{I-6} (citado como~(T)).
  \item[c)] La \emph{teoría de haces}, cuyas referencias principales serán~(G) y~(T); esta teoría proporciona el lenguaje indispensable para interpretar en términos «geométricos» las nociones esenciales del álgebra conmutativa y para «globalizarlas».
  \item[d)] Por último, será útil que el lector esté familiarizado en cierta medida con el \emph{lenguaje funtorial}, que se empleará constantemente en este tratado y para el cual podrá consultar~(M),~(G) y, sobre todo,~(T); los principios de este lenguaje y los principales resultados de la teoría general de los funtores se expondrán con mayor detalle en una obra que los autores de este tratado están preparando.
\end{enumerate}

\bigskip
\begin{center}*\quad*\quad*\end{center}
\bigskip

% R566: fuente francesa intro-fr.tex:24--68, pp. 5--7.
No corresponde a esta introducción dar una descripción más o menos sumaria del punto de vista de los «esquemas» en geometría algebraica, ni la larga lista de razones que han hecho necesaria su adopción y, en particular, la aceptación sistemática de elementos nilpotentes en los anillos locales de las «variedades» que consideramos (lo que necesariamente relega a un segundo plano la noción de aplicación racional en favor de la de aplicación regular o «morfismo»).
Este tratado pretende precisamente desarrollar de manera sistemática el lenguaje de los «esquemas» y demostrará, esperamos, su necesidad.
Aunque sería fácil hacerlo,

\oldpage[I]{6}
tampoco intentaremos dar aquí una introducción «intuitiva» a las nociones desarrolladas en el capítulo primero.
El lector que desee tener una visión preliminar de las materias de este tratado podrá acudir a la conferencia que A.~Grothendieck pronunció en el Congreso Internacional de Matemáticos de Edimburgo en 1958~\cite{I-7} y a la exposición~\cite{I-8} del mismo autor.
El trabajo~\cite{I-14} (citado como~(FAC)) de J.-P.~Serre también puede considerarse una exposición intermedia entre el punto de vista clásico y el de los esquemas en geometría algebraica; por ello, su lectura puede constituir una excelente preparación para la de nuestros \emph{Éléments}.

\bigskip
\begin{center}*\quad*\quad*\end{center}
\bigskip

A título informativo, damos a continuación el plan general previsto para este tratado, sujeto por lo demás a modificaciones posteriores, sobre todo en lo que respecta a los últimos capítulos:

\medskip
\begin{tabular}{r p{0.78\linewidth}}
  Capítulo primero. & --- El lenguaje de los esquemas.\\
  --- \textbf{II}. & --- Estudio global elemental de algunas clases de morfismos.\\
  --- \textbf{III}. & --- Cohomología de los haces algebraicos coherentes. Aplicaciones.\\
  --- \textbf{IV}. & --- Estudio local de los morfismos.\\
  --- \textbf{V}. & --- Procedimientos elementales de construcción de esquemas.\\
  --- \textbf{VI}. & --- Técnica de descenso. Método general de construcción de esquemas.\\
  --- \textbf{VII}. & --- Esquemas en grupos, espacios fibrados principales.\\
  --- \textbf{VIII}. & --- Estudio diferencial de los espacios fibrados.\\
  --- \textbf{IX}. & --- El grupo fundamental.\\
  --- \textbf{X}. & --- Residuos y dualidad.\\
  --- \textbf{XI}. & --- Teorías de intersección, clases de Chern, teorema de Riemann-Roch.\\
  --- \textbf{XII}. & --- Esquemas abelianos y esquemas de Picard.\\
  --- \textbf{XIII}. & --- Cohomología de Weil.\\
\end{tabular}
\medskip

En principio, todos los capítulos se consideran abiertos y siempre podrán añadírseles posteriormente párrafos suplementarios; estos aparecerán en fascículos separados, para reducir los inconvenientes del modo de publicación adoptado.
Cuando uno de esos párrafos esté previsto o en preparación al publicarse un capítulo, se mencionará en el sumario de dicho capítulo, aunque ciertas prioridades hagan que su publicación efectiva deba producirse bastante después.
Para comodidad del lector, ofrecemos en un «Capítulo~0» diversos complementos de álgebra conmutativa, álgebra homológica y teoría de haces, utilizados a lo largo de los capítulos de este tratado; son más o menos conocidos, pero no ha sido posible dar referencias prácticas para ellos.
Se recomienda al lector acudir al capítulo~0 solo durante la lectura del tratado propiamente dicho y en la medida en que los resultados a los que remitimos no le sean

\oldpage[I]{7}
suficientemente familiares.
Pensamos, además, que de esta manera la lectura del tratado podrá ser, para el principiante, un buen método de familiarizarse con el álgebra conmutativa y el álgebra homológica, cuyo estudio, cuando no va acompañado de aplicaciones tangibles, resulta tedioso e incluso deprimente para bastantes personas.

\bigskip
\begin{center}*\quad*\quad*\end{center}
\bigskip

% R567: fuente francesa intro-fr.tex:70--93, pp. 7--8.
Queda fuera de nuestra competencia ofrecer en esta introducción una reseña histórica, siquiera sumaria, de las nociones y los resultados expuestos.
El texto solo contendrá referencias que se consideren particularmente útiles para su comprensión, y únicamente indicaremos el origen de los resultados más importantes.
Al menos desde el punto de vista formal, los temas tratados en nuestra obra son bastante nuevos; esto explica la escasez de referencias a los padres de la geometría algebraica del siglo~\textsc{xix} y de comienzos del~\textsc{xx}, cuyos trabajos solo conocemos de oídas.
Conviene, no obstante, decir aquí unas palabras sobre las obras que han influido más directamente en los autores y contribuido al desarrollo del punto de vista de los esquemas.
Hay que citar en primer lugar el trabajo fundamental (FAC) de J.-P.~Serre, que ha servido de introducción a la geometría algebraica a más de un joven adepto (entre ellos uno de los autores de este tratado) desalentado por la aridez de los clásicos \emph{Foundations} de A.~Weil~\cite{I-18}.
Allí se demuestra por primera vez que la «topología de Zariski» de una variedad algebraica «abstracta» es perfectamente adecuada para aplicarle ciertas técnicas de la topología algebraica y dar lugar, en particular, a una teoría cohomológica.
Además, la definición de variedad algebraica que allí se da es la que se presta más naturalmente a la extensión de esta noción que desarrollamos aquí\footnote{Como nos ha señalado J.-P.~Serre, conviene observar que la idea de definir la estructura de variedad mediante un haz de anillos se debe a H.~Cartan, quien tomó esta idea como punto de partida de su teoría de los espacios analíticos.
Naturalmente, al igual que en geometría algebraica, sería importante dar cabida en «geometría analítica» a los elementos nilpotentes en los anillos locales de los espacios analíticos.
H.~Grauert~\cite{I-5} ha abordado recientemente esta extensión de la definición de H.~Cartan y J.-P.~Serre, y cabe esperar que pronto aparezca una exposición sistemática de geometría analítica en este marco general.
Por otra parte, es evidente que las nociones y las técnicas desarrolladas en este tratado conservan su sentido en geometría analítica, aunque haya que esperar dificultades técnicas mayores en esta última teoría.
Puede preverse que la geometría algebraica, por la sencillez de sus métodos, podrá servir como una especie de modelo formal para desarrollos futuros de la teoría de los espacios analíticos.}.
El propio Serre había observado, además, que la teoría cohomológica de las variedades algebraicas afines podía transcribirse sin dificultad sustituyendo las álgebras afines sobre un cuerpo por anillos conmutativos cualesquiera.
Los capítulos I y II de este tratado y los dos primeros párrafos del capítulo III pueden considerarse, por tanto, esencialmente como transposiciones sencillas, a este marco ampliado, de los resultados principales de (FAC) y de un artículo posterior del mismo autor~\cite{I-15}.
También hemos sacado gran provecho del \emph{Séminaire de Géométrie algébrique} de C.~Chevalley~\cite{I-1}; en particular, el uso sistemático de los «conjuntos construibles» que introdujo se ha revelado muy útil en la teoría de esquemas (véase el cap.~IV).
Le hemos tomado también el estudio de los morfismos desde el punto

\oldpage[I]{8}
de vista de la dimensión (cap.~IV), que se transcribe sin cambios notables al marco de los esquemas.
Conviene señalar, por otra parte, que la noción de «esquemas de anillos locales», introducida por Chevalley, se presta naturalmente a una extensión de la geometría algebraica (sin tener, no obstante, toda la flexibilidad y la generalidad que pretendemos darle aquí); sobre las relaciones entre esta noción y nuestra teoría, véase el capítulo primero, \S~8.
M.~Nagata ha desarrollado tal extensión en una serie de trabajos~\cite{I-9} que contienen numerosos resultados particulares sobre la geometría algebraica sobre anillos de Dedekind\footnote{Entre los trabajos próximos a nuestro punto de vista en geometría algebraica, señalemos el importante trabajo de E.~Kähler~\cite{I-22} y una nota reciente de Chow e Igusa~\cite{I-3}, que retoman, en el marco de la teoría de Nagata-Chevalley, ciertos resultados de (FAC) y dan también una fórmula de Künneth.}.

\bigskip
\begin{center}*\quad*\quad*\end{center}
\bigskip

% R568: fuente francesa intro-fr.tex:95--132, pp. 8--9.
Por último, huelga decir que un libro de geometría algebraica, y sobre todo uno dedicado a sus fundamentos, está necesariamente influido, aunque sea a través de otras personas, por matemáticos como O.~Zariski y A.~Weil.
En particular, la \emph{Théorie des fonctions holomorphes} de Zariski~\cite{I-20}, convenientemente flexibilizada mediante métodos cohomológicos y completada con un teorema de existencia (cap.~III, \S\S~4 y 5), es, junto con la técnica de descenso expuesta en el cap.~VI, una de las principales herramientas empleadas en este tratado, y nos parece una de las más potentes de que dispone la geometría algebraica.

La técnica general en la que se inserta puede esbozarse del siguiente modo (se dará un ejemplo típico en el cap.~IX, al estudiar el grupo fundamental).
Se tiene un morfismo propio (cap.~II) $f:X\to Y$ de una variedad algebraica en otra (más generalmente, de un esquema en otro), que se quiere estudiar cerca de un punto $y\in Y$ para resolver un problema P relativo a una vecindad de $y$.
Se procede por etapas sucesivas:
\begin{enumerate}
  \item[$1^\circ$] Puede suponerse que $Y$ es afín, de modo que $X$ pasa a ser un esquema definido sobre el anillo afín $A$ de $Y$; incluso puede sustituirse $A$ por el anillo local de $y$.
  Esta reducción siempre es fácil en la práctica (cap.~V) y nos lleva al caso en que $A$ es un anillo \emph{local}.
  \item[$2^\circ$] Se estudia el problema planteado cuando $A$ es un anillo local \emph{artiniano}.
  Para que siga teniendo efectivamente sentido cuando no se supone que $A$ sea un dominio, a veces hay que reformular el problema P; de este modo, a menudo se obtiene una mejor comprensión del problema, de carácter «infinitesimal» en esta etapa.
  \item[$3^\circ$] La teoría de esquemas formales (cap.~III, \S\S~3, 4 y 5) permite pasar del caso de un anillo artiniano al de un \emph{anillo local completo}.
  \item[$4^\circ$] Por último, si $A$ es un anillo local cualquiera, la consideración de «secciones multiformes» sobre esquemas adecuados sobre $X$, que aproximan una sección «formal» dada (cap.~IV), permitirá a menudo pasar de un resultado conocido para el esquema

  \oldpage[I]{9}
  deducido de $X$ por extensión de escalares al completado de $A$, a un resultado análogo para una extensión finita bastante sencilla (por ejemplo, no ramificada) de $A$.
\end{enumerate}

Este esbozo muestra la importancia del estudio sistemático de los esquemas definidos sobre un anillo artiniano $A$.
El punto de vista de Serre en su formulación de la teoría local de cuerpos de clases y algunos trabajos recientes de Greenberg parecen sugerir que ese estudio podría emprenderse asociando funtorialmente a tal esquema $X$ un esquema $X'$ sobre el cuerpo residual $k$ de $A$ (suponiendo perfecto ese cuerpo), cuya dimensión sea igual, en los casos favorables, a $n\dim X$, donde $n$ es la longitud de $A$.

En cuanto a la influencia de A.~Weil, baste decir que la necesidad de desarrollar las herramientas para formular con toda la generalidad deseada la definición de «cohomología de Weil» y abordar la demostración\footnote{Para evitar cualquier malentendido, precisemos que esta tarea apenas se ha iniciado cuando se escribe esta introducción y, por tanto, todavía no ha dado lugar a una demostración de las conjeturas de Weil.} de todas las propiedades formales necesarias para establecer sus célebres conjeturas en geometría diofántica~\cite{I-19} ha sido una de las principales motivaciones para redactar este tratado, al igual que el deseo de encontrar el marco natural de las nociones y los métodos habituales de la geometría algebraica y dar a los autores la ocasión de comprender tales nociones y técnicas.

\bigskip
\begin{center}*\quad*\quad*\end{center}
\bigskip

Para terminar, creemos útil advertir a los lectores que, al igual que los propios autores, sin duda encontrarán ciertas dificultades antes de acostumbrarse al lenguaje de los esquemas y convencerse de que las construcciones habituales sugeridas por la intuición geométrica pueden transcribirse en ese lenguaje, esencialmente de una sola manera razonable.
Como en muchas partes de la matemática moderna, la intuición inicial parece alejarse cada vez más del lenguaje apropiado para expresarla con toda la precisión y generalidad deseadas.
En este caso, la dificultad psicológica consiste en la necesidad de trasladar a los objetos de una categoría ya bastante distinta de la categoría de conjuntos (a saber, la categoría de preesquemas o la de preesquemas sobre un preesquema dado) nociones familiares para los conjuntos: productos cartesianos, leyes de grupo, de anillo y de módulo, fibrados, fibrados principales homogéneos, etc.
En el futuro, sin duda será difícil para el matemático eludir este nuevo esfuerzo de abstracción, quizá bastante pequeño, en definitiva, comparado con el que realizaron nuestros predecesores al familiarizarse con la teoría de conjuntos.

\bigskip
\begin{center}*\quad*\quad*\end{center}
\bigskip

Las referencias se darán siguiendo el sistema decimal; por ejemplo, en III, 4.9.3, III indica el capítulo, el número~4 el parágrafo y el~9 la sección de ese parágrafo.
Dentro de un mismo capítulo se omitirá la indicación del capítulo.

